% GENERATED FILE. Edit canonical lessons, then run npm run generate:books.
% canonical-source-hash: fnv1a64:c8c9b18ac03a2d20
% canonical-lessons: ZH-C472-meimao, ZH-C472-lianse, ZH-W472-lian-face, ZH-C472-bianlian, ZH-C472-biaoqing, ZH-C472-etou

\chapter{Faces}
\label{ch:zh-faces}
\hlchaptermodality{472}

\begin{chapteropening}
\textbf{By the end of this chapter:} \emph{I can talk about eyebrows, how someone looks, someone turning nasty, a facial expression and a hot forehead.}

Complete the last lesson of chapter 472: I can talk about eyebrows, how someone looks, someone turning nasty, a facial expression and a hot forehead.
\end{chapteropening}

\section[méimao]{méimao — eyebrow}

\label{lesson:ZH-C472-meimao}



\begin{quote}
Before the new one: say the Mandarin for short-sighted, then the Mandarin for to wink.
\end{quote}

\subsection*{You'll want to know: méimao}
\emph{méimao} — \textquotedblleft{}eyebrow\textquotedblright{}.

\begin{quote}
\emph{tā de méimao bǐ wǒ de hēi} — His eyebrows are darker than mine.
\end{quote}

Its characters come later; for now it is a word for the ear.

\begin{cousinweb}
\emph{méimao} is built from \emph{méi} 'the brow' and \emph{máo} 'hair', both inherited from Middle Chinese, \emph{mij} and \emph{maw} (Guangyun, in Baxter's transcription), which glosses the first, after an older book, as the hair above the eye; the second is said light.
\end{cousinweb}

\begin{grammarlens}[title={What you've built}]
One more word for saying how it was.
\end{grammarlens}

\subsection*{Your turn}
\begin{itemize}
\raggedright
  \item \emph{Say it:} \emph{méimao}
  \item \emph{Say it:} \emph{méimao}, once more
  \item \emph{Say it:} \emph{tā de méimao bǐ wǒ de hēi} — His eyebrows are darker than mine.
  \item \emph{Recall:} say the Mandarin for short-sighted, then the Mandarin for to wink, then say \emph{méimao} again
\end{itemize}

\begin{tcolorbox}[breakable,colback=teal!4,colframe=teal!35!black,title={Before you move on}]
What does méimao mean? (\textquotedblleft{}Eyebrow\textquotedblright{}.) Say it once more.
\end{tcolorbox}
\section[\texorpdfstring{\zh{脸色} (liǎnsè)}{脸色 (liǎnsè)}]{\zh{脸色} (liǎnsè) — complexion; the colour of one's face}

\label{lesson:ZH-C472-lianse}



\begin{quote}
Before the new one: say the Mandarin for to wink, then the Mandarin for an eyebrow.
\end{quote}

\subsection*{You'll want to know: \zh{脸色}}
\textbf{\zh{脸色}} — \emph{liǎnsè} — \textquotedblleft{}complexion; the colour of one's face\textquotedblright{}.

\begin{quote}
\textbf{\zh{你脸色不好是不是病了}} \emph{nǐ liǎnsè bù hǎo shì bu shì bìng le} — You don't look well. Are you ill?
\end{quote}

The new character, in pieces: \zh{脸} — eleven strokes in two pieces: on the left \zh{月} \emph{yuè}, four strokes; on the right seven, as on the right of \zh{检} \emph{jiǎn}: a long stroke falling to the left and one falling to the right, a short bar, two dots, a short stroke falling to the left and a long bar. The next lesson puts it in your hand.

\begin{cousinweb}
\textbf{\zh{脸色}} (\emph{liǎnsè}) is built from \zh{脸} 'face' and \zh{色} 'colour', both inherited from Middle Chinese, \emph{leamX} and \emph{srik} (Guangyun, in Baxter's transcription). The dictionary lists the first under two readings, \emph{tshjem} and \emph{leamX}, and glosses \emph{leamX} as a kind of stew; neither gloss is the face, so how it came to mean face is not settled. It glosses the second 'colour of the face'.
\end{cousinweb}

\begin{grammarlens}[title={What you've built}]
One more word for saying how it was.
\end{grammarlens}

\subsection*{Your turn}
\begin{itemize}
\raggedright
  \item \emph{Say it:} \emph{liǎnsè}
  \item \emph{Say it:} \emph{liǎnsè}, once more
  \item \emph{Say it:} \emph{nǐ liǎnsè bù hǎo shì bu shì bìng le} — You don't look well. Are you ill?
  \item \emph{Recall:} say the Mandarin for to wink, then the Mandarin for an eyebrow, then say \emph{liǎnsè} again
\end{itemize}

\begin{tcolorbox}[breakable,colback=teal!4,colframe=teal!35!black,title={Before you move on}]
What does \zh{脸色} mean? (\textquotedblleft{}Complexion; the colour of one's face\textquotedblright{}.) Say it once more.
\end{tcolorbox}
\section[\texorpdfstring{\zh{脸} (liǎn)}{脸 (liǎn)}]{\zh{脸} — \zh{月}, then the right piece of \zh{检}}

\label{lesson:ZH-W472-lian-face}



\begin{quote}
The word first, then the character inside it.

\begin{itemize}
\raggedright
  \item \emph{Recall:} say \emph{liǎnsè} — \textbf{R1}, one lesson back, before it decays
\end{itemize}

\begin{quote}
\zh{检} \zh{色}
\end{quote}

\emph{Read it:} each one aloud, with nothing said first — \emph{jiǎn}, \emph{sè}
\end{quote}

\subsection*{Script — \zh{月}, then the right piece of \zh{检}}
\hlblockfigure{\detokenize{figures/ZH-W472-lian-face-filmstrip.pdf}}{How \zh{脸} is written, stroke by stroke}

\begin{quote}
\zh{脸}
\end{quote}

\textbf{Two pieces, left before right.} Count the pieces, not the strokes.

\textbf{The left piece, \zh{月}, four strokes, narrow.} The stroke down the left that curves to the left at the foot, the stroke that runs right, turns down and hooks, then two short bars inside.

\textbf{The right piece, seven strokes, as on the right of \zh{检}.} A long stroke falling to the left and a long one falling to the right, like a roof. A short bar under them. Two short dots and a short stroke falling to the left. Then a long bar across the foot.

4 + 7 = eleven.

\subsection*{Your turn}
\begin{itemize}
\raggedright
  \item \emph{Write it:} \zh{脸}, all eleven, with the model in front of you
  \item \emph{Cover:} write it again from memory, and count the pieces, not the strokes
\end{itemize}

\begin{tcolorbox}[breakable,colback=teal!4,colframe=teal!35!black,title={Before you move on}]
How many pieces does \zh{脸} have? (\textbf{Two: four strokes and seven}.) Where have you written the right piece before? (\textbf{On the right of \zh{检}}.)

Source: \href{https://raw.githubusercontent.com/chanind/hanzi-writer-data/68d10a4b21150cae5e1ebbd223eed289cf32d90c/data/\%E8\%84\%B8.json}{Hanzi Writer Data \zh{脸}, pinned stroke paths and medians}.
\end{tcolorbox}
\section[\texorpdfstring{\zh{变脸} (biànliǎn)}{变脸 (biànliǎn)}]{\zh{变脸} (biànliǎn) — to turn hostile suddenly; to change one's attitude abruptly}

\label{lesson:ZH-C472-bianlian}



\begin{quote}
Before the new one: say the Mandarin for an eyebrow, then the Mandarin for the colour of one's face.
\end{quote}

\subsection*{You'll want to know: \zh{变脸}}
\textbf{\zh{变脸}} — \emph{biànliǎn} — \textquotedblleft{}to turn hostile suddenly; to change one's attitude abruptly\textquotedblright{}.

\begin{quote}
\textbf{\zh{我没说什么他怎么突然变脸了}} \emph{wǒ méi shuō shénme tā zěnme tūrán biànliǎn le} — I didn't say anything, so why has he suddenly turned nasty?
\end{quote}

\begin{cousinweb}
\textbf{\zh{变脸}} (\emph{biànliǎn}) is built from \zh{变} 'to change' and \zh{脸} 'face', both inherited from Middle Chinese, \emph{pjenH} and \emph{leamX} (Guangyun, in Baxter's transcription), which glosses the first 'to change'. The sense of a sudden change of attitude is not given there.
\end{cousinweb}

\begin{grammarlens}[title={What you've built}]
One more word for saying how it was.
\end{grammarlens}

\subsection*{Your turn}
\begin{itemize}
\raggedright
  \item \emph{Say it:} \emph{biànliǎn}
  \item \emph{Say it:} \emph{biànliǎn}, once more
  \item \emph{Say it:} \emph{wǒ méi shuō shénme tā zěnme tūrán biànliǎn le} — I didn't say anything, so why has he suddenly turned nasty?
  \item \emph{Recall:} say the Mandarin for an eyebrow, then the Mandarin for the colour of one's face, then say \emph{biànliǎn} again
\end{itemize}

\begin{tcolorbox}[breakable,colback=teal!4,colframe=teal!35!black,title={Before you move on}]
What does \zh{变脸} mean? (\textquotedblleft{}To turn hostile suddenly; to change one's attitude abruptly\textquotedblright{}.) Say it once more.
\end{tcolorbox}
\section[biǎoqíng]{biǎoqíng — facial expression; the look on one's face}

\label{lesson:ZH-C472-biaoqing}



\begin{quote}
Before the new one: say the Mandarin for the colour of one's face, then the Mandarin for to turn hostile suddenly.
\end{quote}

\subsection*{You'll want to know: biǎoqíng}
\emph{biǎoqíng} — \textquotedblleft{}facial expression; the look on one's face\textquotedblright{}.

\begin{quote}
\emph{kàn tā de biǎoqíng wǒ gǎnjué tā bú tài gāoxìng} — From the look on his face, I get the feeling he's not very happy.
\end{quote}

Its characters come later; for now it is a word for the ear.

\begin{cousinweb}
\emph{biǎoqíng} is built from \emph{biǎo} 'the outside' and \zh{情} 'feeling', both inherited from Middle Chinese, \emph{pjewX} and \emph{dzjeng} (Guangyun, in Baxter's transcription). The dictionary glosses the first 'to make clear' and, after an older book, as an outer garment: feeling shown on the outside.
\end{cousinweb}

\begin{grammarlens}[title={What you've built}]
One more word for saying how it was.
\end{grammarlens}

\subsection*{Your turn}
\begin{itemize}
\raggedright
  \item \emph{Say it:} \emph{biǎoqíng}
  \item \emph{Say it:} \emph{biǎoqíng}, once more
  \item \emph{Say it:} \emph{kàn tā de biǎoqíng wǒ gǎnjué tā bú tài gāoxìng} — From the look on his face, I get the feeling he's not very happy.
  \item \emph{Recall:} say the Mandarin for the colour of one's face, then the Mandarin for to turn hostile suddenly, then say \emph{biǎoqíng} again
\end{itemize}

\begin{tcolorbox}[breakable,colback=teal!4,colframe=teal!35!black,title={Before you move on}]
What does biǎoqíng mean? (\textquotedblleft{}Facial expression; the look on one's face\textquotedblright{}.) Say it once more.
\end{tcolorbox}
\section[étóu]{étóu — forehead}

\label{lesson:ZH-C472-etou}



\begin{quote}
Before the new one: say the Mandarin for to turn hostile suddenly, then the Mandarin for the look on one's face.
\end{quote}

\subsection*{You'll want to know: étóu}
\emph{étóu} — \textquotedblleft{}forehead\textquotedblright{}.

\begin{quote}
\emph{wǒ de étóu hěn rè kěnéng fāshāo le} — My forehead's hot. I might have a temperature.
\end{quote}

Its characters come later; for now it is a word for the ear.

\begin{cousinweb}
\emph{étóu} is built from \emph{é} 'the forehead' and \zh{头} 'head', both inherited from Middle Chinese, \emph{ngaek} and \emph{duw} (Guangyun, in Baxter's transcription), which explains the first, after an older book, as the brow.
\end{cousinweb}

\begin{grammarlens}[title={What you've built}]
One more word for saying how it was.
\end{grammarlens}

\subsection*{Your turn}
\begin{itemize}
\raggedright
  \item \emph{Say it:} \emph{étóu}
  \item \emph{Say it:} \emph{étóu}, once more
  \item \emph{Say it:} \emph{wǒ de étóu hěn rè kěnéng fāshāo le} — My forehead's hot. I might have a temperature.
  \item \emph{Recall:} say the Mandarin for to turn hostile suddenly, then the Mandarin for the look on one's face, then say \emph{étóu} again
\end{itemize}

\begin{tcolorbox}[breakable,colback=teal!4,colframe=teal!35!black,title={Before you move on}]
What does étóu mean? (\textquotedblleft{}Forehead\textquotedblright{}.) Say it once more.
\end{tcolorbox}
